\chapter{Conclusions and perspectives}

The present thesis work represents an integral effort in the design and control of the EHEKATL BWB aircraft. Throughout this research, a number of fundamental stages were carried out that addressed from aerodynamic design to the implementation of advanced controllers, culminating in flight simulations in the X-PLANE environment. 

During the design phase, a combination of specialized tools, such as XFLR5, Airoil Maker and Plane Maker, was used to design and evaluate the aerodynamic profile of the EHEKATL BWB aircraft. This process provided a detailed understanding of the pressure distribution along the wing surface, as well as the supporting and drag coefficients. The resulting optimization allowed outstanding aerodynamic efficiency, fundamental to flight performance. XFLR5 released a performance for EHEKATL BWB between 15 and 20 with respect to the polar. $C_L/C_D$ as a function of $\alpha$. In X-PLANE performance was achieved between 15 at the lowest speeds (Mach 0.04) and 27 at the highest speeds (Mach 0.085). A great advantage of our methodology is that all the software is open source and free, except for X-PLANE that to release the simulation time from fifteen minutes to unlimited and all the scenarios you have to buy the full version, however, all the work shown here is developed with the free version. 

The development of the six-degree dynamic model of freedom represented a critical milestone in the understanding of the aircraft's behaviour. By considering the longitudinal and lateral-directional dynamics independently, an accurate representation of individual flight modes was achieved. This included decoupling of nonlinear dynamics, which facilitated a detailed and specific study of each flight phase.

The implementation of the Backstepping and Adaptive Backstepping controllers in both longitudinal and lateral-directional domains proved to be highly effective in stabilizing and controlling the EHEKATL BWB aircraft. These controllers exhibited an exceptional ability to adapt to changing conditions, resulting in accurate and stable control during all flight phases for further implementation in flight control systems.

The flight simulations in the X-PLANE environment provided a significant validation of the effectiveness of the controllers developed. The flight missions, which included take-off and ascent, reconnaissance, as well as descent and landing, confirmed the ability of EHEKATL BWB to operate with precision and control in various conditions and scenarios.



This work lays the foundation for a series of future research and developments, which have the potential to contribute significantly to the field of state-of-the-art aviation:

Exploring techniques for optimisation and adjustment of controller parameters can lead to further improvements in stability and performance in new generation aircraft, as in the case of wing-integrated aircraft, under a variety of operating conditions and flight scenarios. This includes adapting controllers to specific situations, such as changes in atmospheric conditions.

The integration of advanced systems, such as state-of-the-art avionics and innovative propulsion systems, can provide an additional boost to the performance and operational capabilities of EHEKATL BWB. Exploring emerging technologies in areas such as electric propulsion, hydrogen cells, or assignment control systems can open new opportunities in the design of future aircraft.

In perspective, this work lays the foundation for future research and developments in the field of wing-integrated aircraft and flight control systems. The successful implementation of advanced control strategies such as Backstepping and Adaptable Backstepping highlights the road to more efficient, stable and secure aircraft in the future of aviation. EHEKATL BWB emerges not only as a technical achievement, but as a symbol of the Mexican potential in the aerospace landscape.



